


\paragraph{Project (197/200 words).}
Part~II of this dissertation presents \BRACE{} (\cref{sec:irie}), a real-time structural health monitoring platform that the California Department of Transportation (Caltrans) now uses as its frontline tool for post-earthquake bridge assessment.
The platform streams accelerometer data from the field to a cloud service, runs \OpenSees{} models, and generates structural health reports for bridges in the California highway network. 
My research proposes an architecture that sends sensor data over a REST API, queues analyses, and processes results into actionable health scores. 
These summaries reach Caltrans operators within seconds of an event, providing a quantitative basis for closure and inspection decisions. 

The tools that have been developed for this platform are documented online and published as open-source software. 
Most notable among these is \Xara{} (\cref{ch:xara}), a C++ runtime framework for implementing interpreters that has been merged into the central \OpenSees{} repository. 
This framework supports the new \Xara{} Python package, which exposes the first truly embeddable interface for \OpenSees{}, as well as backward-compatible reimplementations of the original Tcl and OpenSeesPy interpreters. 

The visualization features of the platform are released through the Python package \Veux{}.  
This library generates 3D models from \OpenSees{} simulations using the modern glTF format, bypassing the need to distribute heavyweight rendering dependencies. 

\paragraph{Merit (398/400 words).}
Re-engineering the \OpenSees{} interpreters was a significant but necessary endeavor to address systemic issues that prevented their use in mission-critical software like \BRACE{}. 
The existing interpreters employ over 30 global variables, defined across over 20 C++ files, which collectively hold the state of one analysis. These are manipulated indiscriminately through the internal \emph{OPS} API, which consists largely of functions taking no arguments (e.g., \texttt{OPS\_GetAnalysisModel()}) that lack a structured programming paradigm. 
This lack of encapsulation or explicit object ownership leaves the existing interpreters difficult to reason about and plagued with undefined behavior.

The user-facing API provided by OpenSeesPy is similarly unstructured. 
Its extensive reliance on implicit command sequencing—where a function like \texttt{load()} depends on a hidden state set by a previous \texttt{pattern()} call—relegates \OpenSees{} to a standalone scripting tool. 
This global state creates a fragile environment for software development; any function call can have unforeseen side effects on any analysis model, leading to silent numerical corruption or non-deterministic crashes that bypass the interpreter's own exception handling. Consequently, any library that embeds the existing OpenSeesPy must warn its users against combining it with other \OpenSees{}-based libraries for fear that the two libraries will unknowingly corrupt each other's state.

\Xara{} addresses these limitations by stripping global singletons, introducing an explicit object ownership model, and imposing strict encapsulation. 
All commands now operate exclusively on user-owned analysis objects, guaranteeing their state cannot be altered by any routine that has not been explicitly passed a reference to the object. This work, validated by a suite of over \textbf{30} unit tests and merged into the central \OpenSees{} repository, transforms \OpenSees{} into a library-grade engine that can be safely embedded as a predictable ``black-box'' dependency.

In May 2024, Caltrans renewed the BRACE2 platform for a second development phase, expanding its coverage to bridges across California, and work is underway to grant access to over \textbf{40} Caltrans engineers. This adoption solidifies \OpenSees{}RT at the core of the state's seismic response protocol, which has processed over \textbf{150} seismic events since its initial deployment.

Furthermore, the stable, encapsulated API exposed through the companion \Xara{} package is a catalyst for new research, eliminating barriers to integration with modern machine learning toolchains like JAX and PyTorch, which rely crucially on explicit state management. This is already being demonstrated by its adoption in novel studies on parameter sensitivities and surrogate modeling, empowering the community to build more complex and reliable research tools.